% Paste this complete file into main.tex and compile with pdfLaTeX.
\documentclass[11pt,letterpaper]{article}

\usepackage[T1]{fontenc}
\usepackage{lmodern}
\usepackage[margin=0.9in]{geometry}
\usepackage{microtype}
\usepackage{array,booktabs,tabularx}
\usepackage{float}
\usepackage[hidelinks]{hyperref}

\newcolumntype{L}[1]{>{\raggedright\arraybackslash}p{#1}}
\newcolumntype{Y}{>{\raggedright\arraybackslash}X}
\setlength{\parindent}{0pt}
\setlength{\parskip}{0.4em}
\renewcommand{\arraystretch}{1.16}

\newcommand{\menuoption}[2]{%
  \textbf{\(\circ\)\enspace #1}\par
  #2\par\smallskip
}
\newcommand{\ecooption}{%
  \menuoption{Eco}{%
    Adjusts reasoning effort to the task. A complex task may receive
    a higher setting than Light.\par
    \textbf{Design goal}\enspace Lower resource use.\par
    \textbf{Verification status}\enspace Electricity and carbon savings
    unmeasured.%
  }%
}
\newcommand{\fixedoptions}{%
  \menuoption{Light}{Keeps the selected low reasoning effort.}%
  \menuoption{Medium}{Keeps the selected medium reasoning effort.}%
  \menuoption{High}{Keeps the selected high reasoning effort.}%
  \menuoption{Extra High}{Keeps a setting above High.}%
  \menuoption{Max}{Keeps the highest setting in this menu.}%
}
\newcommand{\doi}[1]{\href{https://doi.org/#1}{doi:\,#1}}

\title{Eco Mode Menu Design in Large Language Model Services\\[0.4em]
\large A Comparative Interface Analysis and Two Static Menu Designs}
\author{Jonghoon Ahn}
\date{November 8, 2026}

\begin{document}
\maketitle

\begin{abstract}
Large language model (LLM) services present several kinds of control in
their menus, including model selection, reasoning effort, and automatic
selection. This study examines how an Eco mode could communicate its
environmental objective and adjustment rule within those controls. We
reviewed the official documentation for Gemini Apps, ChatGPT Work, and
Claude chat, along with two Gemini menu screenshots. The comparison
distinguished what a control adjusts, how the setting is decided, its
stated purpose, and the evidence available for its claims. We then
examined two static menu designs with the same Eco option and five fixed
effort levels. Both presented the same functional descriptions and
savings verification status before selection. The grouped menu added
headings that identified the adjustment methods, while both layouts made
the shared descriptions available without opening a further menu. Based
on this comparison, we propose displaying the selected policy and
planned effort for the current task separately, and placing the
resource-saving goal beside the status of savings verification.
\end{abstract}

\textbf{Keywords}\enspace Eco mode, large language models, menu design,
comparative interface analysis, environmental information

\section{Introduction}

When Light and Eco appear side by side, their placement can suggest that
they are settings of the same kind. Light can let a user choose a fixed,
low reasoning effort. Eco can delegate that choice to a policy that
adjusts effort to each task. Both can lead to the same setting for one
request, while following different rules for the next. A menu can make
this distinction explicit by showing both what the user selects and
how the effort level is determined.

The Eco label also introduces an environmental objective. A description
of the setting should identify what the system can adjust and what that
adjustment is intended to achieve. A separate statement can communicate
whether savings have been measured. We compare documented controls in
three services and use the comparison to design and inspect two Eco
menu layouts.

\textbf{RQ1.} Where and how do the three services describe what is
selected, who determines the setting, and the stated purpose of model
and reasoning controls?

\textbf{RQ2.} With option labels and descriptions held constant, how do
two Eco menu layouts present the adjustment rule, environmental
objective, and savings verification status before selection?

\section{Related Work}

GPTFootprint combined an Eco Score, environmental imagery, and
estimated usage information in a ChatGPT extension~\cite{graves2025}.
EcoChat Probe provided estimated energy and carbon feedback to 52
participants over four weeks. Its abstract reports increased awareness
and reflection, alongside no statistically significant change in usage
patterns~\cite{ren2026}. These studies describe how environmental
information can be presented during LLM use.

Klesel and Messer examined choices between options framed in terms of
energy use and performance. Both options used the same underlying
model~\cite{klesel2026}. Kim and colleagues evaluated an LLM interface
that presented response latency and carbon information to university
students~\cite{kim2026}. Building on these examples of environmental
information and model choice, this study examines where an option with
an environmental objective fits among existing model and effort
controls.

\section{Materials and Method}

On November 8, 2026, we reviewed the computer help page for Gemini
Apps, the web documentation for ChatGPT Work, and the model-menu help
page for Claude chat~\cite{google2026,openai2026,anthropic2026}. We
selected these documents to compare how services describe manual
selection, presets, and automatic selection. We also examined two
Korean-language Gemini menu screenshots provided on November 5. We
recorded checkmarks as the selections visible in those screenshots.
Account defaults and rollout conditions require separate evidence.

The unit of analysis was an individual control described in the product
documentation. We treated Power as one preset control because it
combines multiple settings. Within Advanced, we recorded model,
reasoning effort, and speed as separate controls. We adapted the
walkthrough practice of documenting interface steps and wording
described by Light and colleagues~\cite{light2018} to the menu level.
For the service comparison, we inspected public documentation and the
two screenshots. For the design comparison, we applied the same two
illustrative tasks to both layouts and examined the information they
displayed. Table~\ref{tab:criteria} summarizes the recording criteria.

\begin{table}[H]
\centering
\caption{Criteria used to examine the controls.}
\label{tab:criteria}
\small
\begin{tabularx}{\linewidth}{@{}L{0.23\linewidth}Y@{}}
\toprule
\textbf{Criterion} & \textbf{What we recorded} \\
\midrule
Adjustment target &
The model, reasoning effort, speed, or Thinking toggle the control changes. \\
Decision mechanism &
Direct selection, selection of a preset that bundles settings, or
delegation to automatic selection. \\
Stated purpose &
The purpose given in the source, such as response speed, depth, usage
limits, cost, or resource savings. \\
Evidence &
Whether a description comes from documentation or a screenshot,
and where verification information about environmental effects appears. \\
\bottomrule
\end{tabularx}
\end{table}

\section{Comparison of Service Menus}

We identified ten controls across the reviewed materials
(Table~\ref{tab:controls}). The purposes in the table reflect how
providers describe the controls. Comparing realized computation or
performance across similarly named settings would require execution
data.

\begin{table}[H]
\centering
\caption{Controls identified in the public materials. G, W, and C denote
Gemini Apps, ChatGPT Work, and Claude chat, respectively.}
\label{tab:controls}
\small
\begin{tabularx}{\linewidth}{@{}L{0.27\linewidth}L{0.32\linewidth}Y@{}}
\toprule
\textbf{Control} & \textbf{Target and decision mechanism}
& \textbf{Documented purpose or description} \\
\midrule
G1 Gemini model &
Direct choice of model &
Models suited to different response speeds and task complexity. \\
G2 Gemini thinking level &
Direct choice of Low, Medium, or High &
Response speed and reasoning depth. \\
G3 Gemini Auto &
Delegates model and thinking-level selection to the system &
A response suited to the request and management of usage limits. \\
\addlinespace
W1 Work Power &
Choice of a preset combining model and reasoning effort &
Faster processing or deeper reasoning. \\
W2 Work Advanced model &
Direct choice of a supported model &
Specifies which model to use. \\
W3 Work Advanced effort &
Direct choice of a supported reasoning effort &
Reasoning depth suited to the task. \\
W4 Work Advanced speed &
Direct choice of a supported speed setting &
Specifies processing speed. \\
\addlinespace
C1 Claude model &
Direct choice of model &
Changes the model for the current chat. \\
C2 Claude Effort &
Direct choice of a supported effort level &
Relates response depth, time, and usage. \\
C3 Claude Thinking &
Direct control through a toggle where supported &
Configures Thinking or Extended. \\
\bottomrule
\end{tabularx}
\end{table}

The Gemini screenshots showed a model list and a separate thinking-level
submenu with descriptions for Low, Medium, and High. In the
documentation, Gemini Auto delegates the choice of model and thinking
level for a request. ChatGPT Work's Power lets the user choose a preset
that combines settings, with individual controls available under
Advanced. Claude presents model, effort, and thinking controls in a
shared model menu, with availability depending on the model. These
examples distinguish the scope of automation from the composition of
a preset. The selection instructions reviewed here did not specify
comparators or measurement conditions for electricity or carbon
savings~\cite{google2026,openai2026,anthropic2026}.

\section{Eco Menu Designs}

The designs assume one model and use the same six options, order, and
descriptions. Both menus open with all options unselected. A fixed
level retains the setting chosen by the user. Eco is a policy for
selecting reasoning effort by task, with lower resource use stated as
a design goal. Table~\ref{tab:menus} shows the two layouts.

\begin{table}[H]
\centering
\caption{Static menu designs with matched options and descriptions.
The grouped layout adds function headings and spacing.}
\label{tab:menus}
\small
\begin{tabularx}{\linewidth}{@{}Y@{\hspace{1.5em}}Y@{}}
\toprule
\textbf{Flat menu} & \textbf{Grouped by function} \\
\midrule
\textbf{Response mode}\par\smallskip
\ecooption
\fixedoptions
&
\textbf{Response mode}\par\smallskip
\textit{Task-adaptive effort}\par\smallskip
\ecooption
\medskip\textit{Fixed effort levels}\par\smallskip
\fixedoptions
\\
\bottomrule
\end{tabularx}
\end{table}

The flat menu places all options under a common title. The grouped
menu places Eco under \emph{Task-adaptive effort} and Light through Max
under \emph{Fixed effort levels}. In both layouts, the goal and
verification status appear within the Eco entry before selection.

\section{Walkthrough of the Two Designs}

We used two tasks to illustrate the difference between a fixed
setting and an adaptive policy. The first organizes the date, location,
and materials from a workshop announcement. The second compares the
schedules, prices, deadlines, and exceptions for two workshops.
Table~\ref{tab:tasks} gives the planned levels for these scenarios.
They are assignments specified for the walkthrough.

\begin{table}[H]
\centering
\caption{Illustrative effort assignments used in both layouts.}
\label{tab:tasks}
\small
\begin{tabularx}{\linewidth}{@{}Y L{0.21\linewidth}L{0.20\linewidth}@{}}
\toprule
\textbf{Task} & \textbf{Light selected} & \textbf{Eco selected} \\
\midrule
Organize the date, location, and materials from a workshop announcement
& Light & Light \\
Compare schedules, prices, deadlines, and exceptions for two workshops
& Light & High \\
\bottomrule
\end{tabularx}
\end{table}

For the first task, selecting Light and selecting Eco lead to the same
planned level. For the second, fixed Light stays at Light and Eco
assigns High. Recording the selected mode alongside the planned level
for the current task preserves the distinction between the two rules,
even when they produce the same setting.

\begin{table}[H]
\centering
\caption{Information displayed in the initial menu state. An additional
step means opening another menu after the initial menu is open.}
\label{tab:visibility}
\small
\begin{tabularx}{\linewidth}{@{}L{0.39\linewidth}YY@{}}
\toprule
\textbf{Information} & \textbf{Flat menu} & \textbf{Grouped menu} \\
\midrule
Eco adjustment rule and possible higher effort &
Eco description & Eco description \\
Fixed Light setting &
Light description & Light description \\
Resource-saving goal &
Goal line in Eco entry & Goal line in Eco entry \\
Verification status for electricity and carbon savings &
Status line in Eco entry & Status line in Eco entry \\
Headings that distinguish the adjustment methods &
Common menu title & Two function headings \\
Additional menu openings to read the shared descriptions &
0 & 0 \\
\bottomrule
\end{tabularx}
\end{table}

The shared descriptions were visible in both initial menu states
(Table~\ref{tab:visibility}). The grouped layout added a category cue
in the headings, which name the two adjustment methods before the
option descriptions. The number of additional menu openings was the
same. These observations describe the presence and placement of
information in the two layouts. Participant evidence would be needed
to determine what readers notice and how they interpret it.

\section{Discussion}

\subsection{Distinguish What Is Adjusted from Who Sets It}

Gemini Auto delegates the model and thinking level. Work Power lets the
user select a combination of settings~\cite{google2026,openai2026}.
This difference motivates stating exactly what an Eco policy can
adjust. The proposed Eco entry specifies reasoning effort. A design
that also changes the model can list model selection as a separate
adjustment target.

\subsection{Show the Selected Mode and the Planned Task Setting}

In the first scenario, both choices lead to a planned level of Light.
A display that shows only that level leaves the decision rule
ambiguous. Showing \emph{Selected mode Eco} alongside \emph{Planned
effort for this task Light} gives the user a basis for knowing whether
Light was directly selected or assigned by the policy. The distinction
also explains how the setting can change on the next task.

\subsection{Place the Environmental Goal Beside Its Verification Status}

The reviewed Gemini and Work documents describe efficiency, speed,
usage limits, and cost in different contexts~\cite{google2026,openai2026}.
The Eco designs place a resource-saving goal alongside a statement
that electricity and carbon savings are unmeasured. This separates
the purpose of the policy from the evidence available for its effects.
When savings have been measured, the interface can report the
comparison setting, the measured quantity, and the applicable
conditions.

\section{Scope and Conclusion}

The service analysis describes the reviewed documents and the supplied
Gemini screenshots at the time of review. Product controls may vary
with the model, account plan, and interface. The design walkthrough
uses static menus and scripted settings to compare information
placement and access steps. User interpretation, selection, and
willingness to accept performance tradeoffs require participant
evidence. Electricity use and carbon emissions require measurements
under deployment conditions.

In these designs, function headings make the adaptive policy and
fixed effort levels visible as separate categories. Both layouts
provide the same descriptions and verification status within the
initial menu. A design review can use the comparison criteria to
check what the mode adjusts, who sets it, what goal it pursues, and
what evidence is available for the stated effects.

\section*{Use of Generative AI}

Generative AI assisted with organizing public materials, reviewing the
menu designs, drafting the manuscript, and preparing the English text.

\begin{thebibliography}{8}

\bibitem{light2018}
B. Light, J. Burgess, and S. Duguay.
The walkthrough method: An approach to the study of apps.
\emph{New Media \& Society}, 20(3), 881--900, 2018.
\doi{10.1177/1461444816675438}.

\bibitem{graves2025}
N. Graves, V. Larrieu, Y. T. Zhang, J. Peng, V. Nagaraj Rao, Y. Liu,
and A. Monroy-Hern\'andez.
GPTFootprint: Increasing Consumer Awareness of the Environmental
Impacts of LLMs.
\emph{CHI EA 2025}, Article 296. ACM, 2025.
\doi{10.1145/3706599.3719708}.

\bibitem{ren2026}
Y. Ren, L. A. Morais, F. Calegario, F. Krupa, F. P. M. Sim\~oes,
N. Ferreira, and G. Panagiotidou.
EcoChat Probe: How Energy Feedback Shapes End Users' Interpretations
of Everyday Generative AI Use.
\emph{FAccT 2026}, 5194--5218. ACM, 2026.
\doi{10.1145/3805689.3812368}.

\bibitem{klesel2026}
M. Klesel and U. Messer.
Good for the Planet, Bad for Me? Intended and Unintended Consequences
of AI Energy Consumption Disclosure.
\emph{CHI 2026}. ACM, 2026.
\doi{10.1145/3772318.3791015}.

\bibitem{kim2026}
H. Kim, X. Chen, A. Cabral, C. K. Lin, U. Gupta, and J. Hester.
When LLMs Slow Down: How Environmental Impacts Mediate University
Students' LLM Usage.
\emph{ICT4S 2026}, 322--334. IEEE, 2026.
\doi{10.23919/ICT4SCPS00102.2026.00038}.

\bibitem{google2026}
Google.
\href{https://support.google.com/gemini/answer/18648722?co=GENIE.Platform\%3DDesktop\&hl=en}
{\emph{About Gemini models}}.
Gemini Apps Help, Computer. Product documentation.
Accessed November 8, 2026.

\bibitem{openai2026}
OpenAI.
\href{https://learn.chatgpt.com/docs/models?surface=web}{\emph{Models}}.
ChatGPT Learn, ChatGPT Work web interface. Product documentation.
Accessed November 8, 2026.

\bibitem{anthropic2026}
Anthropic.
\href{https://support.claude.com/en/articles/8664678-change-the-model-effort-and-thinking-settings}
{\emph{Change the model, effort, and thinking settings}}.
Claude Help Center. Product documentation. Accessed November 8, 2026.

\end{thebibliography}

\end{document}
